% Part: first-order-logic
% Chapter: syntax-and-semantics
% Section: assignments

\documentclass[../../../include/open-logic-section]{subfiles}

\begin{document}

\olfileid{fol}{syn}{ass}

\olsection{گمارش‌های متغیر}

\begin{explain}
گمارش !!{variable}~$s$ برای \emph{هر} متغیر ارزشی فراهم می‌کند---و
بی‌نهایت متغیر وجود دارد. البته این کار ضروری نیست. فقط از آن رو می‌خواهیم
گمارش‌های !!{variable} به همهٔ !!{variable}s ارزش نسبت دهند که کار را
بسیار آسان‌تر می‌کند. ارزش ترم~$t$ و اینکه !!a{formula}~$!A$ در
!!a{structure} نسبت به~$s$ ارضا می‌شود یا نه، تنها به مقادیری بستگی دارد
که~$s$ به !!{variable}sی درون~$t$ و !!{variable}sی آزاد~$!A$ نسبت
می‌دهد. این محتوای دو گزارهٔ بعدی است. برای دقیق‌کردن مفهوم «بستگی‌داشتن»،
نشان می‌دهیم هر دو گمارش متغیری که دربارهٔ همهٔ متغیرهای~$t$ توافق دارند،
ارزشی یکسان به دست می‌دهند؛ و اگر دو گمارش متغیر دربارهٔ همهٔ متغیرهای
آزاد~$!A$ توافق داشته باشند، $!A$ نسبت به یکی ارضا می‌شود اگر و تنها اگر
نسبت به دیگری ارضا شود.
\end{explain}

\begin{prop}\ollabel{prop:valindep}
اگر !!{variable}sی ترم~$t$ در میان $x_1$، \dots،~$x_n$ باشند و برای
$i = 1$، \dots،~$n$ داشته باشیم $s_1(x_i) = s_2(x_i)$، آنگاه
$\Value{t}{M}[s_1] = \Value{t}{M}[s_2]$.
\end{prop}

\begin{proof}
با استقرا بر پیچیدگی~$t$. در حالت پایه، $t$ می‌تواند !!a{constant} یا یکی
از متغیرهای~$x_1$، \dots،~$x_n$ باشد. اگر $t = c$، آنگاه
$\Value{t}{M}[s_1] = \Assign{c}{M} = \Value{t}{M}[s_2]$. اگر
$t = x_i$، بنا بر فرض گزاره $s_1(x_i) = s_2(x_i)$؛ پس
$\Value{t}{M}[s_1] = s_1(x_i) = s_2(x_i) = \Value{t}{M}[s_2]$.

برای گام استقرا، فرض کنید $t = \Atom{f}{t_1, \dots, t_k}$ و ادعا برای
$t_1$، \dots، $t_k$ برقرار باشد. آنگاه
\begin{align*}
  \Value{t}{M}[s_1] & = \Value{\Atom{f}{t_1, \dots, t_k}}{M}[s_1] \\
  & = \Assign{f}{M}(\Value{t_1}{M}[s_1], \dots, \Value{t_k}{M}[s_1]).
\intertext{برای $j = 1$، \dots،~$k$، !!{variable}sی~$t_j$ در میان
    $x_1$، \dots،~$x_n$ هستند. بنا بر فرض استقرا،
    $\Value{t_j}{M}[s_1] = \Value{t_j}{M}[s_2]$. پس}
\Value{t}{M}[s_1] & = \Value{\Atom{f}{t_1, \dots, t_k}}{M}[s_1] \\
  & = \Assign{f}{M}(\Value{t_1}{M}[s_1], \dots, \Value{t_k}{M}[s_1]) \\
  & = \Assign{f}{M}(\Value{t_1}{M}[s_2], \dots, \Value{t_k}{M}[s_2]) \\
  & = \Value{\Atom{f}{t_1, \dots, t_k}}{M}[s_2] = \Value{t}{M}[s_2].
\end{align*}
\end{proof}

\begin{prop}\ollabel{prop:satindep}
اگر !!{variable}sی آزاد در $!A$ در میان $x_1$، \dots،~$x_n$ باشند و
برای $i = 1$، \dots،~$n$ داشته باشیم $s_1(x_i) = s_2(x_i)$، آنگاه
$\Sat{M}{!A}[s_1]$ اگر و تنها اگر $\Sat{M}{!A}[s_2]$.
\end{prop}

\begin{proof}
از استقرا بر پیچیدگی $!A$ استفاده می‌کنیم. در حالت پایه، که $!A$ اتمی
است، $!A$ می‌تواند یکی از صورت‌های زیر باشد:
\iftag{prvTrue}{$\ltrue$،}{}
\iftag{prvFalse}{$\lfalse$،}{}
$\Atom{R}{t_1, \dots, t_k}$ برای محمول $k$جایگاهی $R$ و
ترم‌های $t_1$، \dots،~$t_k$، یا $\eq[t_1][t_2]$ برای ترم‌های $t_1$
و~$t_2$. در دو حالت اخیر فقط جهت پیش‌روی !!{biconditional} را نشان
می‌دهیم، زیرا برهان جهت عکس متقارن است.

\begin{enumerate}
\tagitem{prvTrue}{%
  \indcase{!A}{\ltrue}{هم $\Sat{M}{\indfrm}[s_1]$ و هم
    $\Sat{M}{\indfrm}[s_2]$.}}{}

\tagitem{prvFalse}{%
  \indcase{!A}{\lfalse}{هم $\Sat/{M}{!A}[s_1]$ و هم
    $\Sat/{M}{!A}[s_2]$.}}{}

\item
  \indcase{!A}{\Atom{R}{t_1, \ldots, t_k}}{فرض کنید
    $\Sat{M}{\indfrm}[s_1]$. آنگاه
    \[
    \langle \Value{t_1}{M}[s_1], \ldots, \Value{t_k}{M}[s_1] \rangle
    \in \Assign{R}{M}.
    \]
    برای $i = 1$، \dots،~$k$، بنا بر \olref{prop:valindep} داریم
    $\Value{t_i}{M}[s_1] = \Value{t_i}{M}[s_2]$. پس همچنین
    $\langle \Value{t_1}{M}[s_2], \ldots, \Value{t_k}{M}[s_2] \rangle
    \in \Assign{R}{M}$؛ بنابراین $\Sat{M}{\indfrm}[s_2]$.}
\item
  \indcase{!A}{\eq[t_1][t_2]}{فرض کنید $\Sat{M}{\indfrm}[s_1]$.
    آنگاه $\Value{t_1}{M}[s_1] = \Value{t_2}{M}[s_1]$. پس
    \begin{align*}
      \Value{t_1}{M}[s_2] & = \Value{t_1}{M}[s_1]
      & \text{(بنا بر \olref{prop:valindep})} \\
      & = \Value{t_2}{M}[s_1]
      & \text{(زیرا $\Sat{M}{\eq[t_1][t_2]}[s_1]$)}\\
      &= \Value{t_2}{M}[s_2]
      & \text{(بنا بر \olref{prop:valindep})،}
    \end{align*}
    پس $\Sat{M}{\eq[t_1][t_2]}[s_2]$.}
\end{enumerate}

اکنون فرض کنید برای همهٔ !!{formula}sی $!B$ که از~$!A$ کم‌پیچیده‌ترند،
$\Sat{M}{!B}[s_1]$ اگر و تنها اگر $\Sat{M}{!B}[s_2]$. گام استقرا به
تفکیک حالت‌هایی پیش می‌رود که عملگر اصلی~$!A$ تعیین می‌کند. در هر حالت
فقط جهت پیش‌روی !!{biconditional} را نشان می‌دهیم؛ برهان جهت عکس متقارن
است. در همهٔ حالت‌ها جز حالت‌های سور، فرض استقرا را بر
!!{subformula}sی~$!B$ از~$!A$ اعمال می‌کنیم. متغیرهای آزاد~$!B$ در میان
متغیرهای آزاد~$!A$ هستند. پس اگر $s_1$ و $s_2$ دربارهٔ متغیرهای آزاد~$!A$
توافق داشته باشند، دربارهٔ متغیرهای آزاد~$!B$ نیز توافق دارند و فرض استقرا
بر~$!B$ اعمال می‌شود.

\begin{enumerate}
\tagitem{defNot}{}{%
  \iftag{probNot}{%
    \indcase!{!A}{\lnot !B}{}}{%
    \indcase{!A}{\lnot !B}{اگر $\Sat{M}{\indfrm}[s_1]$، آنگاه
      $\Sat/{M}{!B}[s_1]$؛ پس بنا بر فرض استقرا،
      $\Sat/{M}{!B}[s_2]$، و ازاین‌رو $\Sat{M}{\indfrm}[s_2]$.}}}

\tagitem{defAnd}{}{%
  \iftag{probAnd}{%
    \indcase!{!A}{!B \land !C}{}}{%
    \indcase{!A}{!B \land !C}{اگر $\Sat{M}{\indfrm}[s_1]$، آنگاه
      $\Sat{M}{!B}[s_1]$ و $\Sat{M}{!C}[s_1]$؛ پس بنا بر فرض استقرا،
      $\Sat{M}{!B}[s_2]$ و $\Sat{M}{!C}[s_2]$. بنابراین
      $\Sat{M}{\indfrm}[s_2]$.}}}

\tagitem{defOr}{}{%
  \iftag{probOr}{%
    \indcase!{!A}{!B \lor !C}{}}{%
    \indcase{!A}{!B \lor !C}{اگر $\Sat{M}{\indfrm}[s_1]$، آنگاه
      $\Sat{M}{!B}[s_1]$ یا $\Sat{M}{!C}[s_1]$. بنا بر فرض استقرا،
      $\Sat{M}{!B}[s_2]$ یا $\Sat{M}{!C}[s_2]$؛ پس
      $\Sat{M}{\indfrm}[s_2]$.}}}

\tagitem{defIf}{}{%
  \iftag{probIf}{%
    \indcase!{!A}{!B \lif !C}{}}{%
    \indcase{!A}{!B \lif !C}{اگر $\Sat{M}{\indfrm}[s_1]$، آنگاه
    $\Sat/{M}{!B}[s_1]$ یا $\Sat{M}{!C}[s_1]$. بنا بر فرض استقرا،
    $\Sat/{M}{!B}[s_2]$ یا $\Sat{M}{!C}[s_2]$؛ پس
    $\Sat{M}{\indfrm}[s_2]$.}}}

\tagitem{defIff}{}{%
  \iftag{probIff}{%
    \indcase!{!A}{!B \liff !C}{}}{%
    \indcase{!A}{!B \liff !C}{اگر $\Sat{M}{\indfrm}[s_1]$، آنگاه یا
      $\Sat{M}{!B}[s_1]$ و $\Sat{M}{!C}[s_1]$، یا
      $\Sat/{M}{!B}[s_1]$ و $\Sat/{M}{!C}[s_1]$. بنا بر فرض استقرا، یا
      $\Sat{M}{!B}[s_2]$ و $\Sat{M}{!C}[s_2]$، یا $\Sat/{M}{!B}[s_2]$
      و $\Sat/{M}{!C}[s_2]$. در هر دو حالت،
      $\Sat{M}{\indfrm}[s_2]$.}}}

\tagitem{defEx}{}{%
  \iftag{probEx}{%
    \indcase!{!A}{\lexists[x][!B]}{}}{%
    \indcase{!A}{\lexists[x][!B]}{اگر $\Sat{M}{\indfrm}[s_1]$،
      $m \in \Domain{M}$ای وجود دارد که
      $\Sat{M}{!B}[\Subst{s_1}{m}{x}]$. بگذارید
      $s_1' = \Subst{s_1}{m}{x}$ و $s_2' =
      \Subst{s_2}{m}{x}$. متغیرهای آزاد~$!B$ در میان $x_1$،
      \dots، $x_n$ و $x$ هستند. داریم $s_1'(x_i) = s_2'(x_i)$، زیرا
      $s_1'$ و $s_2'$ به‌ترتیب واریانت‌های $x$ از $s_1$ و~$s_2$ هستند
      و بنا بر فرض $s_1(x_i) = s_2(x_i)$. بنا بر شیوهٔ تعریف $s_1'$
      و~$s_2'$، داریم $s_1'(x) = s_2'(x) = m$. پس فرض استقرا بر $!B$
      و $s_1'$، $s_2'$ اعمال می‌شود و $\Sat{M}{!B}[s_2']$. بنابراین چون
      $s_2' = \Subst{s_2}{m}{x}$، $m \in \Domain{M}$ای وجود دارد که
      $\Sat{M}{!B}[\Subst{s_2}{m}{x}]$؛ پس
      $\Sat{M}{\indfrm}[s_2]$.}}}

\tagitem{defAll}{}{%
  \iftag{probAll}{%
    \indcase!{!A}{\lforall[x][!B]}{}}{%
    \indcase{!A}{\lforall[x][!B]}{اگر $\Sat{M}{\indfrm}[s_1]$، آنگاه
      برای هر $m \in \Domain{M}$، $\Sat{M}{!B}[\Subst{s_1}{m}{x}]$.
      می‌خواهیم نشان دهیم که برای هر $m \in \Domain{M}$، همچنین
      $\Sat{M}{!B}[\Subst{s_2}{m}{x}]$. پس $m \in \Domain{M}$ را
      دلخواه بگیرید و $s_1' = \Subst{s_1}{m}{x}$ و $s_2' =
      \Subst{s_2}{m}{x}$ را در نظر بگیرید. داریم
      $\Sat{M}{!B}[s_1']$. متغیرهای آزاد~$!B$ در میان $x_1$، \dots،
      $x_n$ و $x$ هستند. داریم $s_1'(x_i) = s_2'(x_i)$، زیرا $s_1'$
      و $s_2'$ به‌ترتیب واریانت‌های $x$ از $s_1$ و~$s_2$ هستند و بنا بر
      فرض $s_1(x_i) = s_2(x_i)$. بنا بر شیوهٔ تعریف $s_1'$ و~$s_2'$،
      داریم $s_1'(x) = s_2'(x) = m$. پس فرض استقرا بر~$!B$ و~$s_1'$،
      $s_2'$ اعمال می‌شود و $\Sat{M}{!B}[s_2']$. این استدلال دربارهٔ هر
      $m \in \Domain{M}$ برقرار است؛ یعنی برای همهٔ~$m \in
      \Domain{M}$، $\Sat{M}{!B}[\Subst{s_2}{m}{x}]$، و بنابراین
      $\Sat{M}{\indfrm}[s_2]$.}}}
\end{enumerate}
بنا بر استقرا، هرگاه !!{variable}sی آزاد در $!A$ در میان $x_1$، \dots،
$x_n$ باشند و برای $i = 1$، \dots،~$n$ داشته باشیم
$s_1(x_i)=s_2(x_i)$، نتیجه می‌شود $\Sat{M}{!A}[s_1]$ اگر و تنها اگر
$\Sat{M}{!A}[s_2]$.
\end{proof}

\begin{probtag}{probNot,probOr,probAnd,probIf,probIff,probEx,probAll}
برهان \olref[fol][syn][ass]{prop:satindep} را کامل کنید.
\end{probtag}

\begin{explain}
!!^{sentence}s هیچ متغیر آزادی ندارند؛ پس هر دو گمارش متغیر دربارهٔ همهٔ
متغیرهای آزاد (یعنی صفر متغیر) هر جمله، چیزهای یکسانی نسبت می‌دهند. بنابراین
گزاره‌ای که اکنون اثبات شد می‌گوید ارضاشدن یا نشدن !!a{sentence} در یک
ساختار نسبت به گمارش متغیر، کاملاً مستقل از گمارش است. این واقعیت را ثبت
می‌کنیم. همین واقعیت تعریف بعدیِ ارضای !!a{sentence} در !!a{structure}
(بدون ذکر گمارش متغیر) را موجه می‌سازد.
\end{explain}

\begin{cor}
\ollabel{cor:sat-sentence}
اگر $!A$ یک !!a{sentence} و $s$ گمارش متغیر باشد، آنگاه
$\Sat{M}{!A}[s]$ اگر و تنها اگر برای هر گمارش متغیر~$s'$،
$\Sat{M}{!A}[s']$.
\end{cor}

\begin{proof}
  بگذارید $s'$ گمارش متغیر دلخواهی باشد. چون $!A$ یک !!a{sentence} است،
  متغیر آزاد ندارد؛ پس هر گمارش متغیر~$s'$ بدیهی است که دربارهٔ همهٔ
  متغیرهای آزاد~$!A$ همان چیزهایی را نسبت می‌دهد که~$s$ می‌دهد. بنابراین
  شرط \olref{prop:satindep} برآورده می‌شود و
  $\Sat{M}{!A}[s]$ اگر و تنها اگر $\Sat{M}{!A}[s']$.
\end{proof}

\begin{defn}
\ollabel{defn:satisfaction}
اگر $!A$ یک !!a{sentence} باشد، می‌گوییم !!a{structure}~$\Struct M$
فرمول~$!A$ را \emph{ارضا می‌کند}، $\Sat{M}{!A}$، اگر و تنها اگر برای همهٔ
گمارش‌های متغیر~$s$، $\Sat{M}{!A}[s]$.
\end{defn}

اگر $\Sat{M}{!A}$، به‌سادگی می‌گوییم \emph{$!A$ در~$\Struct{M}$ صادق
است.} مفهوم ارضا به‌طور طبیعی از !!{sentence}s منفرد به مجموعه‌های
!!{sentence}s گسترش می‌یابد.

\begin{defn}
\ollabel{defn:sat}
اگر $\Gamma$ مجموعه‌ای از !!{sentence}s باشد، می‌گوییم
!!a{structure}~$\Struct M$ مجموعهٔ~$\Gamma$ را \emph{ارضا می‌کند}،
$\Sat{M}{\Gamma}$، اگر و تنها اگر برای همهٔ $!A \in \Gamma$،
$\Sat{M}{!A}$.
\end{defn}

\begin{prop}\ollabel{prop:sentence-sat-true}
  بگذارید $\Struct{M}$ یک !!a{structure}، $!A$ یک !!a{sentence} و $s$
  یک گمارش متغیر باشد. $\Sat{M}{!A}$ اگر و تنها اگر
  $\Sat{M}{!A}[s]$.
\end{prop}

\begin{proof}
تمرین.
\end{proof}

\begin{prob}
\olref[fol][syn][ass]{prop:sentence-sat-true} را اثبات کنید.
\end{prob}

\begin{prop}\ollabel{prop:sat-quant}
  فرض کنید $!A(x)$ تنها $x$ را آزاد در بر دارد و $\Struct M$ یک
  !!a{structure} است. آنگاه:
  \begin{tagenumerate}{prvEx,prvAll}
    \tagitem{prvEx}{$\Sat{M}{\lexists[x][!A(x)]}$ اگر و تنها اگر برای
      دست‌کم یک گمارش متغیر~$s$، $\Sat{M}{!A(x)}[s]$.}{}
    \tagitem{prvAll}{$\Sat{M}{\lforall[x][!A(x)]}$ اگر و تنها اگر برای
  همهٔ گمارش‌های متغیر~$s$، $\Sat{M}{!A(x)}[s]$.}{}
\end{tagenumerate}
\end{prop}

\begin{proof}
تمرین.
\end{proof}

\begin{prob}
\olref[fol][syn][ass]{prop:sat-quant} را اثبات کنید.
\end{prob}

\begin{prob}
\DeclareRobustCommand{\VDash}{\mathrel{||}\joinrel\Relbar}
فرض کنید $\Lang L$ زبانی بدون !!{function}s باشد. با داشتن
!!{structure}~$\Struct M$، !!a{constant} چون $c$ و
$a \in \Domain M$، $\Struct M[a/c]$ را !!{structure}ی تعریف کنید که
درست مانند~$\Struct M$ است، جز آنکه $\Assign{c}{M[a/c]} = a$. برای
!!{sentence}sی~$!A$، $\Struct M \VDash !A$ را چنین تعریف کنید:
\begin{enumerate}
\tagitem{prvFalse}{%
  \indcase{!A}{\lfalse}{نه $\Struct{M} \VDash \indfrm$.}}{}

\tagitem{prvTrue}{%
  \indcase{!A}{\ltrue}{$\Struct{M} \VDash \indfrm$.}}{}

\item \indcase{!A}{\Atom{R}{d_1, \dots, d_n}}{$\Struct M \VDash \indfrm$
  اگر و تنها اگر $\langle \Assign{d_1}{M}, \dots, \Assign{d_n}{M} \rangle \in
  \Assign{R}{M}$.}

\item \indcase{!A}{\eq[d_1][d_2]}{$\Struct M \VDash \indfrm$ اگر و تنها
  اگر $\Assign{d_1}{M} = \Assign{d_2}{M}$.}

\tagitem{prvNot}{%
  \indcase{!A}{\lnot !B}{$\Struct{M} \VDash \indfrm$ اگر و تنها اگر
    نه $\Struct{M} \VDash {!B}$.}}{}

\tagitem{prvAnd}{%
  \indcase{!A}{(!B \land !C)}{$\Struct{M} \VDash
    \indfrm$ اگر و تنها اگر $\Struct{M} \VDash!B$ و
    $\Struct{M} \VDash !C$.}}{}

\tagitem{prvOr}{%
  \indcase{!A}{(!B \lor !C)}{$\Struct{M} \VDash \indfrm$ اگر و تنها اگر
    $\Struct{M} \VDash !B$ یا $\Struct{M} \VDash !C$ (یا هر دو).}}{}

\tagitem{prvIf}{%
  \indcase{!A}{(!B \lif !C)}{$\Struct{M} \VDash \indfrm$ اگر و تنها اگر
    نه $\Struct {M} \VDash !B$ یا $\Struct M \VDash !C$ (یا هر دو).}}{}

\tagitem{prvIff}{%
  \indcase{!A}{(!B \liff !C)}{$\Struct{M} \VDash {\indfrm}$ اگر و تنها
    اگر یا هم $\Struct{M} \VDash {!B}$ و هم
    $\Struct{M} \VDash {!C}$، یا نه $\Struct{M} \VDash {!B}$ و نه
    $\Struct{M} \VDash {!C}$.}}{}

\tagitem{prvAll}{%
  \indcase{!A}{\lforall[x][!B]}{$\Struct{M} \VDash {\indfrm}$ اگر و تنها
    اگر برای همهٔ $a \in \Domain{M}$،
    $\Struct{M[a/c]} \VDash \Subst{!B}{c}{x}$، به شرط آنکه $c$ در~$!B$
    رخ ندهد.}}{}

\tagitem{prvEx}{%
  \indcase{!A}{\lexists[x][!B]}{$\Struct{M} \VDash
  {\indfrm}$ اگر و تنها اگر $a \in \Domain M$ای وجود داشته باشد که
  $\Struct{M[a/c]} \VDash \Subst{!B}{c}{x}$، به شرط آنکه $c$ در~$!B$
  رخ ندهد.}}{}
\end{enumerate}
بگذارید $x_1$، \dots، $x_n$ همهٔ !!{variable}sی آزاد در~$!A$،
$c_1$، \dots، $c_n$ نمادهای ثابتِ ناموجود در~$!A$،
$a_1$، \dots، $a_n \in \Domain M$، و $s(x_i) = a_i$ باشند.

نشان دهید $\Sat{M}{!A}[s]$ اگر و تنها اگر
$\Struct M[a_1/c_1,\dots,a_n/c_n]
\VDash \Subst{\Subst{!A}{c_1}{x_1}\dots}{c_n}{x_n}$.

(این مسئله نشان می‌دهد می‌توان معناشناسی‌ای برای منطق مرتبهٔ اول ارائه کرد
که بدون گمارش متغیر کار کند.)
\end{prob}

\begin{prob}
فرض کنید $f$ نماد تابعی باشد که در~$!A(x,y)$ نیست. نشان دهید
!!a{structure}~$\Struct{M}$ای وجود دارد که
$\Sat{M}{\lforall[x][\lexists[y][!A(x,y)]]}$ اگر و تنها اگر
~$\Struct M'$ای وجود داشته باشد که
$\Sat{M'}{\lforall[x][!A(x,f(x))]}$.

(این مسئله حالت خاصی از نتیجه‌ای است که قضیهٔ اسکولم نامیده می‌شود؛
$\lforall[x][!A(x,f(x))]$ یک \emph{صورت نرمال اسکولم} برای
$\lforall[x][\lexists[y][!A(x,y)]]$ نامیده می‌شود.)
\end{prob}

\end{document}
